% !TeX root = ../drafts/07-instruction-tables.tex
\section{The instruction tables}\label{sec:tables}

Each instruction has its own table.

Write $\rdf{\sep,\addr,\cnt,\val_0,\val_1, \dots}$ for the pair \pull\ $\tup{\sep,\addr,\cnt,\val_0,\val_1, \dots }$, \push\ $\tup{\sep,\addr,\gen\cdot\cnt,\val_0,\val_1, \dots}$: a read of a lookup array (\S\ref{sec:lookup}).

Write $\acc{k}\tup{a;\,v_{\mathrm{old}}\to v_{\mathrm{new}}}$ for the memory access at address $a$ in slot $k$ of a row with clock $\ts$ (\S\ref{sec:memchan}), which is three flushes and one constraint over the access's own columns $\tsprev_k,\glo_k,\ghi_k$ and the counts $r_{\mathrm{lo},k},r_{\mathrm{hi},k}$ of its two range reads:
\begin{align*}
  &\pull\ \tup{\dsep{MEM},a,\tsprev_k,v_{\mathrm{old}}},\quad \push\ \tup{\dsep{MEM},a,\gen^{k}\cdot\ts,v_{\mathrm{new}}},\\
  &\rdf{\dsep{RLO},\glo_k,r_{\mathrm{lo},k}},\qquad \rdf{\dsep{RHI},\ghi_k,r_{\mathrm{hi},k}},\\
  &\tsprev_k\cdot\glo_k=\gen^{k}\cdot\ts\cdot\ghi_k .
\end{align*}
A read is written $\acc{k}\tup{a;\,v}$, for $v_{\mathrm{old}}=v_{\mathrm{new}}=v$. Every table therefore has one constraint per memory access, and the columns listed below leave the accesses' five each implicit. The state flush advances the clock by the stride, $4$ or $20$ (\S\ref{sec:state}).

\subsection{\texttt{XOR64}}\label{sec:tab-xor}

\texttt{XOR64}\ $[o_A,o_B,o_C]$ writes $\loc{o_C}\gets\loc{o_A}+\loc{o_B}$, the $64$-bit field sum (in $\K$).
\begin{itemize}
\item \textbf{Columns:} $\pc,\fp,\ts$; operands $o_A,o_B,o_C$; the two read values $v_A,v_B$; the destination's old value $v_C^{\mathrm{old}}$; bytecode count $r_{\mathrm{bc}}$.
\item \textbf{Constraints:} the three accesses'.
\item \textbf{Flushes:}
  \begin{itemize}
  \item \textbf{State:} \pull\ $\tup{\dsep{ST},\pc,\fp,\ts}$, \push\ $\tup{\dsep{ST},\gen\cdot\pc,\fp,\gen^{4}\cdot\ts}$.
  \item \textbf{Bytecode:} $\rdf{\dsep{BC},\pc,r_{\mathrm{bc}},\opc{XOR64},o_A,o_B,o_C,0,0}$.
  \item \textbf{Memory:} $\acc{0}\tup{\fp\cdot o_A;\,v_A}$, $\acc{1}\tup{\fp\cdot o_B;\,v_B}$, $\acc{2}\tup{\fp\cdot o_C;\,v_C^{\mathrm{old}}\to v_A+v_B}$.
  \end{itemize}
\end{itemize}

\subsection{\texttt{MUL64}}\label{sec:tab-mul}

\texttt{MUL64}\ $[o_A,o_B,o_C]$ writes $\loc{o_C}\gets\loc{o_A} \cdot \loc{o_B}$, the $64$-bit field multiplication (in $\K$).
\begin{itemize}
\item \textbf{Columns:} $\pc,\fp,\ts$; operands $o_A,o_B,o_C$; the two read values $v_A,v_B$; the destination's old value $v_C^{\mathrm{old}}$; bytecode count $r_{\mathrm{bc}}$.
\item \textbf{Constraints:} the three accesses'.
\item \textbf{Flushes:}
  \begin{itemize}
  \item \textbf{State:} \pull\ $\tup{\dsep{ST},\pc,\fp,\ts}$, \push\ $\tup{\dsep{ST},\gen\cdot\pc,\fp,\gen^{4}\cdot\ts}$.
  \item \textbf{Bytecode:} $\rdf{\dsep{BC},\pc,r_{\mathrm{bc}},\opc{MUL64},o_A,o_B,o_C,0,0}$.
  \item \textbf{Memory} (the destination is written with $v_Av_B$, of degree two): $\acc{0}\tup{\fp\cdot o_A;\,v_A}$, $\acc{1}\tup{\fp\cdot o_B;\,v_B}$, $\acc{2}\tup{\fp\cdot o_C;\,v_C^{\mathrm{old}}\to v_Av_B}$.
  \end{itemize}
\end{itemize}

\subsection{\texttt{SET\_CONSTANT}}\label{sec:tab-set}

\texttt{SET\_CONSTANT}\ $[o,k]$ writes $\loc{o}\gets k$.
\begin{itemize}
\item \textbf{Columns:} $\pc,\fp,\ts$; operand $o$; immediate $k$; the cell's old value $v^{\mathrm{old}}$; bytecode count $r_{\mathrm{bc}}$.
\item \textbf{Constraints:} the access's.
\item \textbf{Flushes:}
  \begin{itemize}
  \item \textbf{State:} \pull\ $\tup{\dsep{ST},\pc,\fp,\ts}$, \push\ $\tup{\dsep{ST},\gen\cdot\pc,\fp,\gen^{4}\cdot\ts}$.
  \item \textbf{Bytecode:} $\rdf{\dsep{BC},\pc,r_{\mathrm{bc}},\opc{SET},o,k,0,0,0}$.
  \item \textbf{Memory:} $\acc{0}\tup{\fp\cdot o;\,v^{\mathrm{old}}\to k}$.
  \end{itemize}
\end{itemize}

\subsection{\texttt{DEREF}}\label{sec:tab-deref}

\texttt{DEREF}\ $[o_1,o_2,o_3;\,\dmode]$ reads a pointer and writes the cell it points to. The row accesses three cells:
\begin{itemize}
\item the \emph{pointer} $p$, read at $\fp\cdot o_1$;
\item the \emph{local} value $v_3$, read at $\fp\cdot o_3$;
\item the \emph{target}, written at the dereferenced address $p\cdot o_2$, which held $v_2^{\mathrm{old}}$.
\end{itemize}
The store writes the source chosen by the mode $\dmode$ (\S\ref{sec:vm}): the local cell $v_3$, the return address $\gen^{2}\cdot\pc$, or the frame pointer $\fp$. The mode is two boolean flags $(f_{\pc},f_{\fp})$: $(0,0)$ for $\texttt{deref\_cell}[o_3]$, $(1,0)$ for $\texttt{deref\_pc}$, $(0,1)$ for $\texttt{deref\_fp}$.
\begin{itemize}
\item \textbf{Columns:} $\pc,\fp,\ts$; operands $o_1,o_2,o_3$; flags $f_{\pc},f_{\fp}$; pointer $p$; the local value $v_3$; the target's old value $v_2^{\mathrm{old}}$; bytecode count $r_{\mathrm{bc}}$.
\item \textbf{Constraints:} the three accesses'.
\item \textbf{Flushes:}
  \begin{itemize}
  \item \textbf{State:} \pull\ $\tup{\dsep{ST},\pc,\fp,\ts}$, \push\ $\tup{\dsep{ST},\gen\cdot\pc,\fp,\gen^{4}\cdot\ts}$.
  \item \textbf{Bytecode:} $\rdf{\dsep{BC},\pc,r_{\mathrm{bc}},\opc{DRF},o_1,o_2,o_3,f_{\pc},f_{\fp}}$.
  \item \textbf{Memory}, the target receiving $v_3$, $\gen^{2}\cdot\pc$ or $\fp$ at the three flag settings:
    \begin{align*}
      &\acc{0}\tup{\fp\cdot o_1;\,p},\qquad \acc{1}\tup{\fp\cdot o_3;\,v_3},\\
      &\acc{2}\tup{p\cdot o_2;\,v_2^{\mathrm{old}}\to\bar f\,v_3+f_{\pc}\,(\gen^{2}\cdot\pc)+f_{\fp}\,\fp},
      \qquad \bar f=1+f_{\pc}+f_{\fp}.
    \end{align*}
  \end{itemize}
\end{itemize}

\subsection{\texttt{JUMP}}\label{sec:tab-jump}

\texttt{JUMP}\ $[o_c,o_d,o_f]$ performs a conditional jump on whether the condition $v_{\mathrm{cond}}=\loc{o_c}$ is nonzero. It transfers to $\pc\gets v_{\pc} = \loc{o_d}$, $\fp\gets v_{\fp} = \loc{o_f}$ when $v_{\mathrm{cond}}\neq0$, and defaults to $\pc\gets\gen\cdot\pc$, $\fp\gets\fp$ otherwise. The branch is taken according to a boolean indicator column $b=[\,v_{\mathrm{cond}}\neq0\,]$, certified by a prover-supplied inverse $w$ and 2 constraints:

\begin{itemize}
\item \textbf{Columns:} $\pc,\fp,\ts$; operands $o_c,o_d,o_f$; the condition $v_{\mathrm{cond}}$, the destination $v_{\pc}$ and the frame $v_{\fp}$; inverse $w$ (honest prover sets $w = v_{\mathrm{cond}}^{-1}$ if $v_{\mathrm{cond}}\neq0$, otherwise $w=0$); boolean indicator $b$; bytecode count $r_{\mathrm{bc}}$.
\item \textbf{Constraints:}  $b=v_{\mathrm{cond}}\cdot w$ and $v_{\mathrm{cond}}\,(b+1)=0$, then the three accesses'.
\item \textbf{Flushes:}
  \begin{itemize}
  \item \textbf{State:} \pull\ $\tup{\dsep{ST},\pc,\fp,\ts}$, \push\ $\tup{\dsep{ST},b\,v_{\pc}+b\,(\gen\cdot\pc)+\gen\cdot\pc,\;b\,v_{\fp}+b\,\fp+\fp,\;\gen^{4}\cdot\ts}$.
  \item \textbf{Bytecode:} $\rdf{\dsep{BC},\pc,r_{\mathrm{bc}},\opc{JMP},o_c,o_d,o_f,0,0}$.
  \item \textbf{Memory:} $\acc{0}\tup{\fp\cdot o_c;\,v_{\mathrm{cond}}}$, $\acc{1}\tup{\fp\cdot o_d;\,v_{\pc}}$, $\acc{2}\tup{\fp\cdot o_f;\,v_{\fp}}$.
  \end{itemize}
\end{itemize}

\subsection{\texttt{BLAKE2S}}\label{sec:tab-blake2s}

\texttt{BLAKE2S}\ $[o_{m_0},o_{m_1},o_{m_2},o_{m_3},o_{\mathit{cv}},o_{\mathit{out}},o_{\md}]$ compresses a $64$-byte message block under a $32$-byte chaining value into a $32$-byte result, with the BLAKE2s byte counter and flags in the two metadata cells from $\fp\cdot o_{\md}$ (\S\ref{sec:vm}). Every cell holds one $8$-byte word: message chunk $i$ spans the cells $\gen^{k}\cdot\fp\cdot o_{m_i}$, $k<2$, the four chunks addressed independently; the chaining value spans $\gen^{k}\cdot\fp\cdot o_{\mathit{cv}}$ and the result $\gen^{k}\cdot\fp\cdot o_{\mathit{out}}$, $k<4$; the metadata spans $\gen^{k}\cdot\fp\cdot o_{\md}$, $k<2$.

Flock~\cite{Flock26} proves the compression relation with a batched R1CS over a Boolean witness (Annex~\ref{annex:flock}). The witness is packed into $\qflock$ via Ring-Switching (Annex~\ref{annex:rs}), with $64$ bits per $\K$-element.

The message, result, chaining value and metadata are all committed in $\qflock$: the eighteen words $m_{i,k}$ ($i<4$, $k<2$), $\mathit{out}_k$ and $\mathit{cv}_k$ ($k<4$), and $\md_k$ ($k<2$), one $\qflock$ slot each. The generalized R1CS matrices take the chaining value, byte counter, and the two finalization flags as free input: the memory interactions via the bus bind all four.
\begin{itemize}
\item \textbf{Columns:} $\pc,\fp,\ts$; operands $o_{m_0},o_{m_1},o_{m_2},o_{m_3},o_{\mathit{cv}},o_{\mathit{out}},o_{\md}$; the four result cells' old values $\mathit{out}^{\mathrm{old}}_k$; bytecode count $r_{\mathrm{bc}}$. The eighteen words are committed in $\qflock$, no duplicate columns needed. The eighteen addresses are the products $\gen^{k}\cdot\fp\cdot o$.
\item \textbf{Constraints:} the eighteen accesses'. The compression relation is handled by Flock.
\item \textbf{Flushes:}
  \begin{itemize}
  \item \textbf{State:} \pull\ $\tup{\dsep{ST},\pc,\fp,\ts}$, \push\ $\tup{\dsep{ST},\gen\cdot\pc,\fp,\gen^{20}\cdot\ts}$: eighteen slots do not fit a stride of $4$.
  \item \textbf{Bytecode:} $\rdf{\dsep{BC},\pc,r_{\mathrm{bc}},\opc{B2S},o_{m_0},o_{m_1},o_{m_2},o_{m_3},o_{\mathit{cv}},o_{\mathit{out}},o_{\md}}$.
  \item \textbf{Memory}, the fourteen reads taking slots $0$ to $13$ and the four writes of the result the last four, so that the result may land on cells the compression read:
    \begin{align*}
      &\acc{2i+k}\tup{\gen^{k}\cdot\fp\cdot o_{m_i};\,m_{i,k}}\quad(i<4,\ k<2),\\
      &\acc{8+k}\tup{\gen^{k}\cdot\fp\cdot o_{\mathit{cv}};\,\mathit{cv}_{k}}\quad(k<4),\\
      &\acc{12+k}\tup{\gen^{k}\cdot\fp\cdot o_{\md};\,\md_{k}}\quad(k<2),\\
      &\acc{14+k}\tup{\gen^{k}\cdot\fp\cdot o_{\mathit{out}};\,\mathit{out}^{\mathrm{old}}_{k}\to\mathit{out}_{k}}\quad(k<4).
    \end{align*}
  \end{itemize}
\end{itemize}
